%!TEX root = ../../thesis.tex

\section{Where the Approach Helps}\label{sec:discussion-advantages}

% TODO: Interpretation, not a re-listing of \cref{chap:results}. Every claim here must
% point at a specific result, and claims the results do not support must not appear.
%
% Candidate arguments, each to be kept only if the evidence carries it:
%   - deep, sequentially dependent constraints: where a coverage-guided climber stalls
%     because no single mutation improves coverage, a learned value function that
%     carries credit across steps does not (\cref{sec:results-input-actions}),
%   - sample efficiency where an execution is expensive: under emulation, spending
%     compute per decision to spend fewer executions is the right trade, and that is
%     the regime this domain sits in (\cref{sec:rl-sample-efficiency}),
%   - grammar discovery from raw coverage: the agent converges on deeply nested
%     structure on both recursive-descent parsers without being told the grammar
%     (\cref{sec:results-corpus-actions}, \cref{tab:best-inputs}),
%   - a planned mutation is worth more than a random one at the same execution budget,
%     which is what the iteration-bounded control arm isolates
%     (\cref{sec:target-cjson}).
%
% State the scope of each claim: which target, which baseline, which currency. An
% advantage demonstrated only on \cref{sec:target-sequence}, where the reward was
% designed, is a claim about credit assignment, not about fuzzing --- keep that
% distinction visible here rather than only in \cref{sec:discussion-threats}.
%
% The bound belongs here as much as the advantage: picohttpparser
% (\cref{sec:target-picohttpparser}) is the target where the approach buys nothing,
% and naming the property it lacks --- coverage that grows with input structure --- is
% what turns three coverage curves into a design lesson.
